\scanpage{30}

Si on dispose d'une théorie suffisamment souple des types d'homotopie, on devrait pouvoir en déduire le type d'homotopie de $X$ comme ``somme amalgamée'' des types d'homotopie des trois \deleted{topos}\inserted{``espaces''} $\unsure{\mathcal V_{Y,X}^*|V_{Y,X}^*}$, $\unsure{\mathcal V_{Y,X}|V_{Y,X}}$ et $X^*$ (dont les deux premiers sont de nature ``infinitésimale''). Par exemple, on doit avoir pour les \unsure{espaces} déduits en ``tuant les $\pi_i$ pour $i\geq2$'' -- donc, essentiellement, pour les groupoïdes fondamentaux, un ``théorème de van Kampen'', qui s'exprime par le fait que le diagramme des groupoïdes fondamentaux

(8)
\diagram{\begin{tikzcd}[ampersand replacement=\&,column sep=small]
\Pi_1\unsure{\mathcal V_{Y,X}^*|V_{Y,X}^*}\arrow[rr,"p"]\arrow[d,"i"']\&\&\Pi_1\unsure{\mathcal V_{Y,X}|V_{Y,X}}\ (\xleftarrow{\sim}\Pi_1Y)\arrow[d,dotted]\\
\Pi_1X^*\&\bigl(\arrow[r,dotted]\&\Pi_1X\ \bigr)
\end{tikzcd}}

\margin{29}
